\documentclass[11pt,a4paper]{article}
\usepackage{gtu-el-yazisi}
\begin{document}
\notbilgi{Alan Teorisi}{Dielektrik ortam}{Ders13.pdf}

\sayfa{1}
Boş olmayan ortam

\elcizimi[0.75]{sekiller/Ders13-s1-1.png}[1]
Belli bir zaman sonra dengeye erişip bir iç yük oluşturacak

Boş uzayda

\elcizimi[0.3]{sekiller/Ders13-s1-2.png}[1]
\[ \vec E = -\mathrm{grad}\,V \qquad \vec D = \varepsilon_0\vec E \qquad \vec D\parallel\vec E \]
\[ \Delta V = -\frac{\rho}{\varepsilon_0} = \frac{-\rho_{\text{dış}}}{\varepsilon_0} \]
\[ \mathrm{div}\,\vec D = \mathrm{div}(\varepsilon_0\vec E) = \rho = \rho_{\text{dış}} \]

(Ortamın fiziksel varlığından ötürü dış yük tarafından yaratılan $\vec E$'nin etkisiyle oluşan yüklerin yoğunluğu)

$\rho = \rho_{\text{dış}} + \bar\rho$ \qquad artık ortamı boş gibi düşünebiliriz

\elcizimi[0.3]{sekiller/Ders13-s1-3.png}[1]
\[ \mathrm{div}(\varepsilon_0\vec E) = \rho = \rho_{\text{dış}} + \bar\rho \]
\[ \bar\rho = -\mathrm{div}\Bigl(\varepsilon_0\,\overbrace{\chi(\vec E)}^{\text{? kapa katsayısı ortamın duyarlılığıdır}}\,\vec E\Bigr) \]

\begin{align*}
\mathrm{div}(\varepsilon_0\vec E) - \bar\rho &= \rho_{\text{dış}} \\
\intertext{div lineerdir}
\mathrm{div}(\varepsilon_0\vec E) + \mathrm{div}\bigl(\varepsilon_0\chi(\vec E)\,\vec E\bigr) &= \rho_{\text{dış}} \\
\mathrm{div}\Bigl(\varepsilon_0\,\underbrace{\bigl(1+\chi(\vec E)\bigr)}_{\varepsilon_r>1}\cdot\vec E\Bigr) &= \rho_{\text{dış}}
\end{align*}

İçinde bir takım dipollerin oluştuğu ortama dielektrik ortam denir. Böyle bir ortamda $\vec P = \varepsilon_0\,\chi(\vec E)\cdot\vec E$, $\vec E$ ile aynı yöndedir, dolayısıyla da
\[ \underbrace{\vec P \parallel \vec E}_{\text{aynı yönde ve paraleldir.}} \;\Rightarrow\; \chi(\vec E) > 0 \]

\[ \mathrm{div}\bigl(\underbrace{\varepsilon_0\varepsilon_r}_{\varepsilon}\cdot\vec E\bigr) = \rho_{\text{dış}} \;\Rightarrow\; \mathrm{div}(\varepsilon\cdot\vec E) = \rho_{\text{dış}} \]

\[ \vec D = \varepsilon\vec E \]
\formulkutu{\mathrm{div}\,\vec D = \rho_{\text{dış}}}
ortam basitse kapa elektrik alana bağımlı değil sabittir.

Basit ortamda $\varepsilon = \varepsilon_0\cdot\varepsilon_r \;\Rightarrow\; \varepsilon > \varepsilon_0$ \qquad $\varepsilon_r = \dfrac{\varepsilon}{\varepsilon_0} > 1$ bağıl dielektrik sabiti

$\Delta V = -\dfrac{\rho_{\text{dış}}}{\varepsilon}$ \quad ortam basitse

\noindent
\begin{minipage}[t]{0.5\linewidth}
Boşluk: $\mathrm{div}\,\vec D = \rho_{\text{dış}}$ \quad ($\vec D \to \varepsilon_0\vec E$)

$\mathrm{div}\,\vec E = -\rho_{\text{dış}}/\varepsilon_0$ \qquad $\mathrm{div}(\vec D) = \rho_{\text{dış}}$ \quad ($\vec D \to \varepsilon\vec E$)

$\Delta V = -\rho_{\text{dış}}/\varepsilon_0 \;\longrightarrow\; V = \dfrac{q}{4\pi\varepsilon_0 r}$
\end{minipage}\hfill
\begin{minipage}[t]{0.46\linewidth}
$\mathrm{rot}\,\vec E = 0$

$\mathrm{div}\,\vec D = \rho_{\text{dış}} \;\Rightarrow\; \mathrm{div}(\varepsilon\vec E) = \rho_{\text{dış}}$

$\vec F = q\,\vec E$ \qquad $\mathrm{div}\,\vec E = \dfrac{\rho_{\text{dış}}}{\varepsilon}$
\end{minipage}

\sayfa{2}
\elcizimi[0.35]{sekiller/Ders13-s2-1.png}[2]
$q=1$ ise \quad $V = \dfrac{1}{4\pi\varepsilon\cdot r}$

\[ V = \frac{1}{4\pi\varepsilon}\int\frac{\rho}{r'}\,dv \]
\[ \vec E = -\mathrm{grad}\,V \qquad \vec D = \varepsilon\vec E \qquad \mathrm{rot}\,\vec E = 0 \]

\medskip\noindent\hrulefill\medskip

\elcizimi[0.3]{sekiller/Ders13-s2-2.png}[2]
\[ \vec n = \frac{d\vec S}{dS} \]

\[ \mathrm{rot}\,\vec E=0 \;\Rightarrow\; \{\mathrm{rot}\,\vec E\} + \vec n\times(\vec E_2-\vec E_1)\big|_S\,\delta(s) = \{0\} + 0\cdot\delta(s) \]

\[ \left\{\begin{array}{l}\mathrm{rot}\,\vec E_1 : v_1\\ \mathrm{rot}\,\vec E_2 : v_2\end{array}\right\} + \vec n\times\bigl(\vec E_2\big|_S-\vec E_1\big|_S\bigr)\,\delta(s) \]
\[ \overset{\makebox[6cm][r]{$0$}}{\rule{6cm}{0.4pt}} \]
\[ = \left\{\begin{array}{ll}0 & v_1\\ 0 & v_2\end{array}\right\} + 0\cdot\delta(s) \]

\[ \vec n\times\bigl(\vec E_2\big|_S-\vec E_1\big|_S\bigr) = 0 \]
\[ \Rightarrow\; \vec n\times\vec E_2\big|_S = \vec n\times\vec E_1\big|_S \;\Rightarrow\; \vec E_1^{\,t}\big|_S = \vec E_2^{\,t}\big|_S \]

\elcizimi[0.35]{sekiller/Ders13-s2-3.png}[2]
\begin{align*}
\vec n\times\vec A &= \vec n\times\bigl(A^t\,\vec t + A^n\,\vec n\bigr) \\
 &= A^t\,\underbrace{(\vec n\times\vec t)}_{\vec t\,'} + A^n\,\underbrace{\vec n\times\vec n}_{0} \\
\vec n\cdot\vec A &= \vec n\cdot\bigl(A^t\,\vec t + A^n\,\vec n\bigr) \\
 &= A^t\,\vec n\cdot\vec t + A^n\,\vec n\cdot\vec n = A^n
\end{align*}

\sayfa{3}
\[ \mathrm{div}\,\vec D = \rho \]
\[ \left\{\begin{array}{l}\mathrm{div}\,\vec D_1 : v_1\\ \mathrm{div}\,\vec D_2 : v_2\end{array}\right\} + \vec n\,(\vec D_2-\vec D_1)\big|_S\text{\footnotemark} = \left\{\begin{array}{l}\rho_1 ; v_1\\ \rho_2 ; v_2\end{array}\right\} + \rho_s\,\delta(s) \]
\footnotetext{Claude notu: muhtemelen sol tarafta da $\delta(s)$ çarpanı yazılmalıydı.}

\begin{align*}
\mathrm{div}\,\vec D_1 &= \rho_1 & \vec n\,(\vec D_2-\vec D_1)\big|_S &= \rho_s \\
\mathrm{div}\,\vec D_2 &= \rho_2 & \Rightarrow\ \vec n\,\vec D_2\big|_S - \vec n\,\vec D_1\big|_S &= \rho_s
\end{align*}

\[ D_2^n\big|_S - D_1^n\big|_S = \rho_s \;\;=\; -\varepsilon_2\frac{\partial V_2}{\partial n}\Big|_S + \varepsilon_1\frac{\partial V_1}{\partial n}\Big|_S = \rho_s \]
\[ \Rightarrow\; +\varepsilon_2\frac{\partial V_2}{\partial n}\Big|_S - \varepsilon_1\frac{\partial V_1}{\partial n}\Big|_S = -\rho_s \]

\[ \vec E = -\mathrm{grad}\,V = \left\{\begin{array}{l}-\mathrm{grad}\,V_1 : v_1^{\text{iç}}\\ -\mathrm{grad}\,V_2 : v_2^{\text{dış}}\end{array}\right\} + \overbrace{\vec n\,(V_2-V_1)\big|_S}^{0}\,\delta(s) \]

\formulkutu{V_1\big|_S = V_2\big|_S}

\begin{align*}
\Delta V_1 &= -\frac{\rho_1}{\varepsilon_1} && v_1\text{'de} \\
\Delta V_2 &= -\frac{\rho_2}{\varepsilon_2} && v_2\text{'de}
\end{align*}

\elcizimi[0.45]{sekiller/Ders13-s3-1.png}[3]
\begin{gather*}
\vec E = -\mathrm{grad}\,V = -\mathrm{grad}\,V_s = 0 \\
\vec D = \varepsilon\vec E = 0 \\
\mathrm{div}\,\vec D = \rho = 0
\end{gather*}

\[ \varepsilon_2\frac{\partial V_2}{\partial n} = -\rho_s \]
\[ V_1\big|_S = V_s\big|_S = V_s = \text{sabit} \]

\formulkutu{\begin{aligned} V_1\big|_S &= V_s\big|_S = V_s = \text{sabit} \\ V_2\big|_{s=\infty} &< \infty \end{aligned}}

\sayfa{4}
Yüzey üzerinde potansiyel sabitse ve içeride yük yoksa iç bölgede potansiyel sabittir, elektrik alan de sabittir, ve içerde $\rho=0$ \,/\!/

\elcizimi[0.6]{sekiller/Ders13-s4-1.png}[4]
\begin{gather*}
V_2 = V_d = \text{sabit} \\
\vec D_2 = 0 \\
\rho_2 = 0
\end{gather*}

\[ \mathrm{div}\,\vec D = \rho \;\Rightarrow\; \Delta V_1 = \frac{\partial^2 V_1}{\partial x^2} + \frac{\partial^2 V_1}{\partial x^2}\text{\footnotemark} + \frac{\partial^2 V_1}{\partial z^2} = 0 \]
\footnotetext{Claude notu: muhtemelen ikinci terim $\partial^2 V_1/\partial y^2$ olmalı.}

\elcizimi[0.35]{sekiller/Ders13-s4-2.png}[4]
\begin{align*}
\frac{d^2V_1}{dx^2} &= 0 \;\Rightarrow\; V_1 = Ax + B \\
V_1\Big|_{x=0} &= A\cdot 0 + B = V_0 \;\Rightarrow\; B = V_0 \\[2mm]
V_1\Big|_{x=d} &= Ad + V_0 = V_d
\end{align*}

\[ \rho_{s1} = \vec e_x\Bigl(\vec D_2\Big|^{0}_{x=d} - \vec D_1\Big|_{x=d}\Bigr) \]
\[ \rho_{s1} = -\vec D_1\Big|_{x=d}\,\vec e_x = -\varepsilon_1\,\vec E_1\,\Big|\,\vec e_x \qquad x>d \]
\[ \rho_s^{0} = -\rho_s^{d} \]

\end{document}
